\chapter{Conclusion}
\label{ch:conclusion}

QLens shows that a person who owns only a messy spreadsheet can be taken, step by step and with every
step explained, to a verified third-normal-form database. Type inference, key search and dependency
discovery run inside PostgreSQL through PL/pgSQL and dynamic SQL; near-duplicates are found with
\code{pg\_trgm} over a GIN index; repairs are transactional and audited; the synthesis is Bernstein's
algorithm with a key relation, property-tested for 3NF and losslessness; and the apply step proves on
the real rows, with \code{EXCEPT ALL} in both directions inside one transaction, that nothing was lost.
On the fee register every planted problem is found and the expected four-table design rebuilds all
3{,}000 rows exactly. Around this sits a 22-table metadata schema that uses generated columns, \code{CHECK}
rules, partial unique indexes, triggers and window-function views, and a per-workspace login role that
isolates users by privilege. 229 backend and 63 frontend tests pass.

\paragraph{Lessons.} Dependencies in real data are statistical before they are logical, so measuring
violations ($g_3$) and letting the user decide was the central design decision. Normalization needs clean
data first: a rule broken by eight typos would otherwise be lost and the design would keep a redundant
column. Set-oriented checks belong in the database. Constraints are documentation that executes: they
caught our own mistakes, such as an applied-design rule that let \code{NULL} pass until it was written as
\code{IS TRUE}. And losslessness, a theorem in the textbook, becomes in practice a query.

\paragraph{Limitations.} Left-hand sides are limited to two columns and keys to three; noise filtering
relies on thresholds, so the user keeps the final say; one flat sheet of up to 100{,}000 rows is
normalized at a time, and only to 3NF; trigram similarity cannot know that two people share a name; and
the evidence comes from synthetic samples and tests, not a user study. The query editor, design canvas
and assistant are \stwip{}, and the performance laboratory and step-through visualizer are \stplan.

\paragraph{Future work.} Completing the sandboxed editor and the canvas with its validated assistant; a
performance laboratory that renders \code{EXPLAIN (ANALYZE, BUFFERS)} as a tree and times index
experiments on a scaled copy; a clause-by-clause visualizer for teaching; wider dependency search with
sampling and multivalued dependencies for 4NF; and resource limits for hosted, multi-tenant use, which the
per-workspace role design was chosen to make possible.
